\section{Additional parameters and numerical verification details}
\label{app:extended_checks}

This appendix collects the parameters, transfer-matrix comparisons and mesh diagnostics. We report the algebraic checks separately from comparisons with published band diagrams, for which the available data are incomplete.

\subsection{Parameter summary}
\label{app:parameter_summary}
Table~\ref{tab:s1_parameters} gives the inputs for the main calculations and bilayer comparisons.
% Parameter summary split into three conventional LaTeX tables for reliable elsarticle output.
\begin{table}[p]
\centering
\caption{Case-H homogeneous model: material parameters and computational settings}
\label{tab:s1_parameters}
\footnotesize
\renewcommand{\arraystretch}{1.18}
\begin{tabular}{@{}p{0.11\textwidth}p{0.25\textwidth}p{0.56\textwidth}@{}}
\toprule
Study & Parameter group & Values used \\
\midrule
Case H & Cell and base material & $L=1.0\,\mathrm{m}$; $\lambda=\mu=1.0\,\mathrm{Pa}$; $\rho=1.0\,\mathrm{kg/m^3}$ \\
Case H & Gradient and inertia lengths & $l_{\mathrm{iso}}=0.20\,\mathrm{m}$; $\ell_{\mathrm{i}}=0.20\,\mathrm{m}$ ($\ell_{\mathrm{i}}^2=0.04\,\mathrm{m^2}$) \\
Case H & Anisotropy scaling & Volume-equivalent axes $l_1=l_{\mathrm{iso}}\sqrt{\mathrm{AR}}$ and $l_2=l_{\mathrm{iso}}/\sqrt{\mathrm{AR}}$; $\mathrm{AR}\in\{1,2,3,5,7,10\}$ \\
Case H & Design-map sampling & $\mathrm{AR}\in\{1,2,3,5,7,10\}$ and $\theta\in\{0^\circ,15^\circ,30^\circ,45^\circ,60^\circ,75^\circ,90^\circ\}$; 42 sampled $(\theta,\mathrm{AR})$ points \\
Case H & Frequency-sweep sampling & $\mathrm{AR}=5$: $\theta=0^\circ$--$90^\circ$ in $2.5^\circ$ steps; $\theta=45^\circ$: $\mathrm{AR}\in\{1,2,3,5,7,10\}$ \\
Case H & Steering sampling & $\bar{k}=0.5$; $\mathrm{AR}\in\{5,10\}$; $\theta=0^\circ,15^\circ,\ldots,90^\circ$; ring direction sampled at $5^\circ$ \\
Case H & Band paths and branches & $\Gamma$--$X$--$M$--$\Gamma$, 40 intervals per leg (121 path nodes); four lowest physical branches \\
Case H & Production resolution & $16\times16$ BFS elements per cell ($h/L=1/16$); homogeneous results use the listed $16\times16$ mesh unless stated otherwise \\
Case H & Zone scans & $41\times81$ half-zone grid for complete-gap/IFC evaluations; the Case-C quoted gap uses the separate grid stated below \\
Case H & Mesh-convergence point & $\bmk=(0.31\pi/L,0.22\pi/L)$; $n\times n$, $n\in\{4,8,16,32\}$ \\
Case H & Numerical check tolerances & Hermiticity $<10^{-12}$; symmetry discrepancies $<10^{-10}$; acoustic-slope discrepancy $<10^{-6}$; energy-flux residual $<10^{-6}$ \\
\bottomrule
\end{tabular}
\end{table}

\begin{table}[p]
\centering
\caption{Case-C periodic composite: material parameters and computational settings}
\label{tab:s1_caseC_parameters}
\footnotesize
\renewcommand{\arraystretch}{1.18}
\begin{tabular}{@{}p{0.11\textwidth}p{0.25\textwidth}p{0.56\textwidth}@{}}
\toprule
Study & Parameter group & Values used \\
\midrule
Case C & Epoxy matrix & $\rho_m=1142\,\mathrm{kg/m^3}$; $\mu_m=1.48\,\mathrm{GPa}$; $\lambda_m=4.57\,\mathrm{GPa}$ \cite{zhanwei2010} \\
Case C & YBCO inclusion & $\rho_i=6333\,\mathrm{kg/m^3}$; $\mu_i=37.0\,\mathrm{GPa}$; $\lambda_i=95.0\,\mathrm{GPa}$ \cite{zhanwei2010} \\
Case C & Geometry and gradient lengths & $r_0/L=0.30$ (area fraction $28.27\%$); $\ell_m/L=0.10$; $\ell_i/L=0.20$ \\
Case C & Frequency scale and grid & $c_{t,m}=1138.4\,\mathrm{m/s}$; $\omega_{0,C}=3576.4\,\mathrm{rad/s}$; quoted gap on $4\times4$ mesh and $21\times21$ Brillouin-zone grid \\
\bottomrule
\end{tabular}
\end{table}

\begin{table}[p]
\centering
\caption{ Compare with cases from the transfer matrix: Material properties and model settings.}
\label{tab:s1_transfer_parameters}
\footnotesize
\renewcommand{\arraystretch}{1.18}
\begin{tabular}{@{}p{0.11\textwidth}p{0.25\textwidth}p{0.56\textwidth}@{}}
\toprule
Study & Parameter group & Values used \\
\midrule
B1 & AlN layer A & $a_A=0.01\,\mathrm{m}$; $\rho_A=3230\,\mathrm{kg/m^3}$; $c_{33,A}=3.9\times10^{11}\,\mathrm{Pa}$; $a_{3,A}=8.4\times10^{-11}$ in source units \cite{li2024anchorB} \\
B1 & BaTiO$_3$ layer B & $a_B=0.01\,\mathrm{m}$; $\rho_B=5800\,\mathrm{kg/m^3}$; $c'_{33,B}=1.62\times10^{11}\,\mathrm{Pa}$; $a'_{3,B}=1.26\times10^{-8}$ in source units \cite{li2024anchorB} \\
B1 & Classical setting and normalisation & $b=a_A+a_B=0.02\,\mathrm{m}$; $l=l_1=f=0$; $\omega_0=2.2422\times10^6\,\mathrm{rad/s}$ \\
B2 & Published gradient parameters & Printed values $l=10^{-5}$, $l_1=2\times10^{-5}$, $f=F=0$, and layer ratios $L=L_1=5$; the units and dimensional versus normalised interpretation of the lengths are unspecified in the source \cite{li2024anchorB} \\
B2 & Interpretations evaluated & Dimensional micro-scale ($a=l=10^{-5}\,\mathrm{m}$); dimensional 1-cm layer ($a=0.01\,\mathrm{m}$, $l=10^{-5}\,\mathrm{m}$); and normalised-length reading ($\bar l=10^{-5}$). These give different curves; B2 is not externally validated \\
B3 & Pb/brass material and thickness & $a_1=a_2=10\,\mu\mathrm{m}$; $\mu_1=2.3\times10^{10}\,\mathrm{Pa}$; $\rho_1=7500\,\mathrm{kg/m^3}$; $\mu_2=1.288\times10^9\,\mathrm{Pa}$; $\rho_2=1177.5\,\mathrm{kg/m^3}$ \cite{li2023anchorA} \\
B3 & Gradient, inertia, and frequency scale & $c_1/a_1^2=0.15$; $d_1/a_1=0.25$; $c_2/c_1=d_2/d_1=1.5$; $\omega_{0,B3}=4.114\times10^8\,\mathrm{rad/s}$; the first-layer ratios are inherited from the paper's comparison example, not reported on the plotted panel \\
\bottomrule
\end{tabular}
\end{table}



\subsection{Transfer-matrix comparisons and validation limits}
\label{app:transfer_checks}
\subsubsection{Classical AlN/BaTiO$_3$ bilayer}
For the classical bilayer, $a_A=a_B=0.01\,\mathrm{m}$ and the period is $b=0.02\,\mathrm{m}$. Its transfer-matrix trace reduces to the Rytov relation in the main manuscript \cite{li2024anchorB,zhengwei2009}. The homogeneous single-layer identity and identical-material bilayer reduction have maximum relative residuals of $6.47\times10^{-16}$ and $7.22\times10^{-16}$. The heterogeneous trace agrees with the Rytov relation to $2.22\times10^{-16}$. The calculated normalised stop-band intervals are $[0.4806,0.5209]$, $[0.9812,1.0215]$, and $[1.4993,1.5028]$. Only a visual comparison with the published plot is possible. Without numerical source curves, a percentage error cannot be calculated reliably.

\subsubsection{Gradient-elastic bilayers and validation limits}
In the AlN/BaTiO$_3$ gradient case, the published plot lists $l=10^{-5}$ and $l_1=2\times10^{-5}$.  The length symbols are not specified to be dimensional or normalised in the surrounding text but this is indicated by the plot. The paper defines $\bar l=l/b$ with $b=a_A+a_B=0.02\,\mathrm{m}$. Reading the printed $10^{-5}$ as $\bar l$ gives $l=2\times10^{-7}\,\mathrm{m}$. Using a single layer width as the scale gives $10^{-7}\,\mathrm{m}$. Together with the dimensional-micro and dimensional-macroscale readings, these choices produce different transfer-matrix curves. For the dimensional-micro interpretation, the calculated stop bands are $[0.078,0.242]$, $[0.378,0.383]$, $[0.742,0.760]$, $[0.785,1.103]$, and $[1.123,1.435]$. The homogeneous-limit and identical-material reductions each have an algebraic residual of $4.10\times10^{-57}$. These identities will only check the transfer-matrix algebra. They don't test agreement with published dispersion curves. No information can be provided on numerical source data, and the parameter convention remains undetermined. \cite{li2024anchorB}.

For the Pb/brass shear-horizontal case, the evaluated ratios are $c_1/a_1^2=0.15$, $d_1/a_1=0.25$, and $c_2/c_1=d_2/d_1=1.5$. Li et al. use the same numerical example throughout to illustrate $c_1$ and $d_1$. They are not mentioned in the cited comparison panel. The calculated stop-band intervals are $[0.340,1.024]$, $[1.423,1.867]$, and $[2.477,2.911]$. The homogeneous-limit identity and identical-material reduction have relative residuals of $1.37\times10^{-13}$ and $4.19\times10^{-14}$, respectively. The printed dipolar layer matrix is reproduced to machine precision. This establishes algebraic consistency only. The cited panel does not provide the parameter set, normalisation or numerical curves needed for external validation \cite{li2023anchorA,liweizhou2016}.

\subsubsection{Pure-bending micro-beam limit}
The pure-bending micro-beam reduction matches the closed-form relations of Papargyri-Beskou and Beskos to better than $10^{-15}$ in both evaluated limits \cite{papargyribeskou2009}. This analytical benchmark is in addition to the plot comparisons involving bilayers.

\subsection{Numerical consistency and mesh-convergence details}
\label{app:numerical_checks}
This appendix gives further details of the eight numerical checks and the Case-H mesh study.


The normalised directional widths for the representative Case-H designs in the main-text gap summary are given in Table~\ref{tab:s2_gap_ratio}. The selected rows are reported to be in agreement between the $8\times8$ and $16\times16$ meshes to the precision selected. There are 126 gap-measure values reported across the three metrics from the design map of 42 configurations ($6$ aspect-ratio values $\times$ $7$ orientations). The maximum mesh-to-mesh changes are $5.5\times10^{-3}$ for $\Delta_{GX}$, $5.8\times10^{-3}$ for $\Delta_{\mathrm{path}}$, and $5.7\times10^{-3}$ for $\Delta_{\mathrm{complete}}$.
\begin{table}[htbp]
\centering
\caption{Normalised directional-gap widths for the representative Case-H designs.}
\label{tab:s2_gap_ratio}
% Alt text (retained in alt_text_for_revision.txt): The normalised directional-gap values are negative in every listed design. More elongated cases show larger negative widths, with orientation dependence appearing for the most elongated configuration.
\begin{tabular*}{0.85\textwidth}{@{\extracolsep{\fill}}c c c@{}}
\toprule
$\mathrm{AR}$ & $\theta$ & $\Delta/\bar{\omega}_{\mathrm{mid}}$ \\
\midrule
$1$  & $0^\circ$  & $-0.09$ \\
$1$  & $45^\circ$ & $-0.09$ \\
$1$  & $90^\circ$ & $-0.09$ \\
$3$  & $0^\circ$  & $-0.13$ \\
$3$  & $45^\circ$ & $-0.13$ \\
$3$  & $90^\circ$ & $-0.13$ \\
$5$  & $0^\circ$  & $-0.16$ \\
$5$  & $45^\circ$ & $-0.16$ \\
$5$  & $90^\circ$ & $-0.16$ \\
$10$ & $0^\circ$  & $-0.23$ \\
$10$ & $45^\circ$ & $-0.15$ \\
$10$ & $90^\circ$ & $-0.23$ \\
\bottomrule
\end{tabular*}
\end{table}

\begin{itemize}
  \item The stiffness and mass Hermiticity residuals are $2.01\times10^{-16}$ and $5.63\times10^{-17}$.The reciprocal-lattice shift affects the reduced matrices at most to $4.0\times10^{-16}$, and the spectrum at most to $9.66\times10^{-16}$
  \item The greatest frequency change across the sampled orientations is $9.63\times10^{-16}$ at $\mathrm{AR}=1$. The spectral discrepancy of the anisotropic rotation--swap comparison is $6.44\times10^{-16}$. If the axis swap is not performed, the result is $0.179$.
  \item The best long-wave speed discrepancy is $1.25\times10^{-8}$. The deviation at $\bar{k}=10^{-2}$, $1.80\times10^{-6}$, may include an expected $O(\bar{k}^2)$ dispersion and it is not considered a discretisation error.
  \item At $\Gamma$, The stiffness operator is expected to have 2 modes of translation, which is indeed correct. The lowest eigenvalue of the interior is $\omega^2=3.028$. The minimum mass-matrix eigenvalue is $9.9\times10^{-5}>0$.
  \item At $\bar{k}=200$, the bounded phase-speed result differs from its asymptote by $2.85\times10^{-4}$. The residual between the energy-flux and the group-velocity is $6.91\times10^{-10}$. It employs exact Case-H harmonic fields, cycle-averaged flux and group velocities for both branches of acoustic wave propagation at various propagation angles.
\end{itemize}

\begin{table}[htbp]
\centering
\caption{Detailed convergence information for expanded mesh: successive changes, eigensolver reproducibility.}
\label{tab:s3_mesh}
% Alt text (retained in alt_text_for_revision.txt): The acoustic-frequency error and successive mesh changes decrease as the Case-H mesh is refined. The finest listed mesh is identified as resolution-limited because its error is below measured eigensolver reproducibility.
\begin{tabular}{@{}c r r r r l@{}}
\toprule
Mesh & $\bar{\omega}_T$ & Relative error & \shortstack{Successive change\\(relative)} & Reproducibility & Use in fit \\
\midrule
$4\times4$ & 1.164855406908 & $1.51\times10^{-8}$ & --- & $2.90\times10^{-14}$ & Included \\
$8\times8$ & 1.164855390213 & $7.59\times10^{-10}$ & $1.43\times10^{-8}$ & $1.06\times10^{-13}$ & Included \\
$16\times16$ & 1.164855389383 & $4.63\times10^{-11}$ & $7.12\times10^{-10}$ & $5.92\times10^{-13}$ & Included \\
$32\times32$ & 1.164855389329 & $2.48\times10^{-15}$ & $4.63\times10^{-11}$ & $1.26\times10^{-12}$ & Resolution-limited \\
\bottomrule
\end{tabular}

\end{table}

Table~\ref{tab:s3_mesh} contains the mesh data. The fit to the $4^2$, $8^2$ and $16^2$ points gives $p=4.17$ (95\% CI $[3.15,5.20]$).The $32^2$ result is not considered resolution-limited as the standard deviation is smaller than three times the measured standard deviation in reproducibility. The operational mesh-change floor $\varepsilon_\Delta=4.63\times10^{-11}$, is different from the eigensolver reproducibility at the finest level.
\subsection{Gap resolution in the context of case and parameter sensitivity}
\label{app:caseC_resolution}
The influence of the finite-element mesh and the Brillouin-zone sampling on the Case-C complete gap is separated in Table~\ref{tab:s4_caseC_mesh}. At the $4\times4$ mesh, the four leg/path values are $\Delta[\Gamma\text{--}X]=2.6114$, $\Delta[X\text{--}M]=2.7563$, $\Delta[M\text{--}\Gamma]=3.2871$, and $\Delta[\mathrm{path}]=2.5732$, compared with $\Delta[\mathrm{complete}]=2.5732$. On the coarsest mesh, the gap at $X$ is larger than the full gap since the lowest frequency of the upper band is at an off $X$ point on the $\Gamma$--$M$ diagonal. The values of $X$ (and the complete gaps) are consistent to the degree shown, beginning with $8\times8$.

\begin{table}[htbp]
\centering
\caption{Case-C: Complete Gap and $X$-point Gap values with mesh refinement.}
\label{tab:s4_caseC_mesh}
% Alt text (retained in alt_text_for_revision.txt): The complete gap for the epoxy--YBCO cell decreases monotonically as the finite-element mesh is refined from $4\times4$ to $64\times64$, indicating that the coarse-mesh gap is not converged.
\begin{tabular}{@{}c r r r r r@{}}
\toprule
Mesh & DOFs & $\Delta_{\mathrm{complete}}$ & $\Delta_X$ & Decrement & Decrement ratio (next/this) \\
\midrule
$4\times4$ & 200 & 2.5732 & 2.7563 & --- & --- \\
$8\times8$ & 648 & 2.2504 & 2.2504 & 0.3228 & 0.552 \\
$16\times16$ & 2312 & 2.0722 & 2.0722 & 0.17814 & 0.554 \\
$32\times32$ & 8712 & 1.9736 & 1.9736 & 0.09861 & 0.565 \\
$64\times64$ & 33800 & 1.9179 & 1.9179 & 0.05575 & --- \\
\bottomrule
\end{tabular}

\end{table}

At the $4\times4$ mesh, increasing the Brillouin-zone grid from $11\times11$ to $21\times21$ leaves the complete gap at $2.5732$. The $41\times41$ grid gives $2.5728$. The upper edge is shifted from $7.1430$ to $7.1426$ when the minimum along $\Gamma$--$M$ is better-resolved. For $8\times8$ and $16\times16$ meshes, the $11\times11$ and $21\times21$ grids agree to four decimals. For $32\times32$ and $64\times64$, they agree to six decimals. With the finite-element mesh fixed at $4\times4$, using $4\times4$, $6\times6$, and $8\times8$ Gauss points per cut element gives $X$-point gaps of $2.7563$, $2.7226$, and $2.7862$, respectively (changes of $-1.22\%$ and $+1.08\%$ relative to the $4\times4$ rule). Varying the radius over $r_0/L=0.20$--$0.40$ covers filling fractions of $12.6\%$--$50.3\%$. The complete gap changes from $0.8288$ ($16.61\%$) at $r_0/L=0.20$ to $2.5732$ ($43.94\%$) at $0.30$, then to $2.4426$ ($37.00\%$) at $0.40$. As seen in Table~\ref{tab:s4_caseC_mesh}, the maximum of the mesh is not converged, with the width gaps decreasing as you get to the higher meshes.

